% !TEX root = ../OzicaneNeuronskeMreze.tex
\section{Formalni neuron i principi neuralnih mreža}
% Profesor nije naveo izvore za ovu sekciju: dodati vlastite
% (npr. McCulloch i Pitts 1943, Rosenblatt 1958, udžbenik o neuralnim mrežama).
% Kratko: ova sekcija samo uvodi pojmove koje ostatak poglavlja koristi.

\subsection{Model formalnog neurona}
% Ulazi, težine, težinska suma, prag, izlaz. Slika neurona nacrtana u TikZ-u.

\subsection{Aktivacijske funkcije}
% Prag (step), sigmoid, ReLU; zašto je bitno koju funkciju hardver može realizovati.

\subsection{Perceptron i pravilo učenja}
% Pravilo korekcije težina; učenje = mijenjanje težina, tj. „programa“ mreže.

\subsection{Višeslojne mreže}
% Slojevi, propagacija unaprijed, kratko o učenju propagacijom greške unazad.

\subsection{Operacija MAC kao osnova hardverske realizacije}
% Množenje i akumulacija (multiply-accumulate): operacija koju ubrzava sav hardver
% iz narednih sekcija. Ovdje je uvesti, kasnije se samo pozivati na nju.
